\documentclass[10pt, aspectratio=169]{beamer}
\setbeameroption{hide notes} 

% Tema y colores UTFSM
\usetheme{Madrid}
\usecolortheme{default}
\definecolor{usmblue}{RGB}{0, 75, 135}
\definecolor{usmgold}{RGB}{200, 160, 50}
\setbeamercolor{palette primary}{bg=usmblue,fg=white}
\setbeamercolor{palette secondary}{bg=usmblue!80!black,fg=white}
\setbeamercolor{palette tertiary}{bg=usmblue!60!black,fg=white}
\setbeamercolor{palette quaternary}{bg=usmblue,fg=white}
\setbeamercolor{structure}{fg=usmblue}
\setbeamercolor{item projected}{bg=usmblue}
\setbeamertemplate{navigation symbols}{}
\setbeamertemplate{caption}[numbered]

% Paquetes
\usepackage[utf8]{inputenc}
\usepackage[spanish]{babel}
\usepackage{graphicx}
\usepackage{amsmath,amssymb}
\usepackage{booktabs}
\usepackage{hyperref}
\usepackage{multirow}
\usepackage{tikz}
\usetikzlibrary{shapes.geometric, arrows.meta, positioning, calc}

% Información de la Portada
\title[Reconstrucción EEG con GSP y TRSS]{Reconstrucción de Señales EEG mediante Procesamiento de Señales en Grafos e Interpolación Temporal}
\author[Carlos Saldivia]{Carlos Saldivia}
\institute[UTFSM]{
  Departamento de Ingeniería Electrónica\\
  Universidad Técnica Federico Santa María\\[1.5ex]
  \textit{Directores: Dr. Alejandro Weinstein, Dr. Matías Zañartu}
}
\date{Defensa de Tesis de Magíster -- 2026}

\begin{document}

\begin{frame}[plain]
    \titlepage
    \note{Buenos días al comité evaluador. Hoy presentaré mi defensa de tesis de magíster, donde abordamos el problema de interpolación de canales EEG perdidos, combinando teoría de grafos y regularización temporal. Este trabajo busca superar las severas limitaciones de los métodos actualmente considerados estándar en la neurociencia computacional.}
\end{frame}

\begin{frame}{Contenidos}
    \tableofcontents[hideallsubsections]
    \note{La presentación está estructurada en 5 bloques. Empezaremos contextualizando la motivación clínica. Luego explicaré el marco teórico de Grafos y regularización temporal. Dedicaré tiempo a la metodología, ya que allí radica un aporte importante para evitar el sesgo estadístico (Optuna). Finalmente revisaremos el cruce de rendimiento, la degradación comparativa visual, el test de ablación de nuestro parámetro clave, y las conclusiones.}
\end{frame}

% --- SECCIÓN 1: INTRODUCCIÓN ---
\section{Introducción y Motivación}

\begin{frame}{El Electroencefalograma (EEG) y la Pérdida de Información}
    \begin{itemize}
        \item El EEG es fundamental para el diagnóstico clínico (ej., epilepsia) y el diseño de Interfaces Cerebro-Computador (BCI).
        \item \textbf{El Problema Real:} La pérdida de información en uno o más canales no suele ser aleatoria; frecuentemente vemos \textit{parches enteros} de electrodos perdidos por secado del gel o desconexión del amplificador.
        \item \textbf{Impacto Analítico:}
        \begin{itemize}
            \item Interrumpe el análisis topográfico y localización de fuentes.
            \item Distorsiona morfologías transitorias críticas, como los ERPs (ej. P300).
            \item Obliga a descartar \textit{trials} enteros.
        \end{itemize}
    \end{itemize}
    
    \vspace{0.2cm}
    \begin{block}{Hipótesis de Trabajo}
        La interpolación puramente espacial falla bajo pérdida contigua severa. Incorporar la \textbf{inercia temporal} neurofisiológica mediante Procesamiento de Señales en Grafos (GSP) permite reconstruir fielmente la señal, preservando sus componentes espectrales y transitorios.
    \end{block}
    \note{Aunque el EEG es estándar, es muy frágil. Lo importante aquí es la pérdida contigua o de parche (nearby). Perder un electrodo aislado es fácil de arreglar, pero perder 4 o 5 vecinos juntos destruye la topografía. Y descartar todo ese segmento es inaceptable cuando recolectar datos es tan costoso.}
\end{frame}

\begin{frame}{Limitaciones del Estado del Arte (Enfoques Clásicos)}
    \textbf{Interpolación Geométrica (Ej: Spherical Splines de Perrin, 1989)}
    \begin{itemize}
        \item Estándar \textit{de facto} en \texttt{EEGLAB} y \texttt{MNE-Python}.
        \item \textit{Supuesto:} Asume que la cabeza es una esfera estática y la señal decae exclusivamente por distancia euclidiana.
        \item \textbf{Fallas Estructurales:}
        \begin{enumerate}
            \item Opera estrictamente en el dominio espacial \textbf{instante por instante} ($t$ fijo).
            \item \textbf{Ignora la inercia neurofisiológica:} el voltaje cortical no "salta" de forma instantánea; tiene derivadas continuas.
        \end{enumerate}
    \end{itemize}
    
    \vspace{0.3cm}
    \begin{block}{Consecuencia}
        Las reconstrucciones euclidianas introducen ruido de alta frecuencia y, bajo degradación severa, colapsan invirtiendo la fase temporal de la señal.
    \end{block}
    \note{Las Spherical Splines son elegantes pero ingenuas. Como operan frame a frame, no tienen 'memoria'. Si hay un parche de electrodos perdidos, la Spline fuerza una superficie 3D que muchas veces hace que un pico positivo termine siendo reconstruido como un valle negativo.}
\end{frame}

% --- SECCIÓN 2: MARCO TEÓRICO ---
\section{Marco Teórico: GSP y TRSS}

\begin{frame}{Del Espacio Euclidiano al Grafo (GSP)}
    \begin{columns}
        \begin{column}{0.5\textwidth}
            \textbf{Modelado GSP}
            \begin{itemize}
                \item Se abandona la esfera y se pasa a un \textbf{Grafo Ponderado} $\mathcal{G} = (\mathcal{V}, \mathcal{E}, \mathbf{W})$.
                \item $W_{ij}$ captura la conectividad funcional o topológica.
                \item El \textbf{Laplaciano} $\mathbf{L} = \mathbf{D} - \mathbf{W}$ penaliza la variación espacial entre vecinos.
            \end{itemize}
        \end{column}
        \begin{column}{0.5\textwidth}
            \begin{block}{Regularización Tikhonov}
                Optimización instantánea:
                \[ \hat{\mathbf{x}} = \arg\min_{\mathbf{z}} \| \mathbf{y} - \mathbf{S} \mathbf{z} \|_2^2 + \alpha (\mathbf{z}^\top \mathbf{L} \mathbf{z}) \]
            \end{block}
            \textit{El parámetro $\alpha$ controla qué tan suave queremos que sea la distribución de voltajes en el grafo.}
        \end{column}
    \end{columns}
    \note{El GSP nos libera de la restricción de la distancia en 3D. El Laplaciano L nos permite decir 'este nodo y este otro son funcionalmente vecinos'. Tikhonov mejora las splines, pero sigue resolviendo el problema instante por instante.}
\end{frame}

\begin{frame}{Regularización Espacio-Temporal (TRSS)}
    El método \textbf{Temporal Regularized Sobolev Smoothness (TRSS)} optimiza la ventana completa simultáneamente $\mathbf{X} \in \mathbb{R}^{N \times T}$:
    
    \vspace{0.2cm}
    \begin{exampleblock}{Función Objetivo TRSS}
    \[ \hat{\mathbf{X}} = \arg\min_{\mathbf{Z}} \underbrace{\| \mathbf{S}\mathbf{X} - \mathbf{S}\mathbf{Z} \|_F^2}_{\text{1. Fidelidad}} + \underbrace{\alpha \, \text{tr}(\mathbf{Z}^\top \mathbf{L} \mathbf{Z})}_{\text{2. Suavidad Espacial}} + \underbrace{\beta \| \mathbf{Z}\mathbf{D}_t \|_F^2}_{\text{3. Suavidad Temporal}} \]
    \end{exampleblock}
    \vspace{0.1cm}
    
    \begin{columns}
        \begin{column}{0.33\textwidth}
            \begin{alertblock}{1. Fidelidad}
                \centering Obliga a que la solución coincida con los electrodos \textbf{sanos} registrados.
            \end{alertblock}
        \end{column}
        \begin{column}{0.33\textwidth}
            \begin{block}{2. Espacio ($\alpha$)}
                \centering Penaliza variaciones bruscas entre electrodos \textbf{vecinos} según el Grafo (Laplaciano $\mathbf{L}$).
            \end{block}
        \end{column}
        \begin{column}{0.33\textwidth}
            \begin{block}{3. Tiempo ($\beta$)}
                \centering \textbf{Inercia}. Penaliza derivadas temporales grandes ($\mathbf{D}_t$). Evita saltos no fisiológicos.
            \end{block}
        \end{column}
    \end{columns}
    \note{Esta ecuación es el núcleo. TRSS resuelve toda una época de una sola vez. Al penalizar Z D_t, impone la inercia neurofisiológica: 'Si vas a cambiar el voltaje drásticamente en el milisegundo siguiente, debes tener evidencia en los canales vecinos, sino, mantente continuo'.}
\end{frame}

% --- SECCIÓN 3: METODOLOGÍA ---
\section{Metodología y Diseño Experimental}

\begin{frame}{Pipeline Metodológico de la Investigación}
    \centering
    \tikzset{
        box/.style={rectangle, rounded corners, draw=usmblue, fill=usmblue!10, thick, minimum height=0.8cm, text centered, font=\small},
        arrow/.style={-{Stealth[scale=1.2]}, thick, usmblue}
    }
    \begin{tikzpicture}[node distance=1.5cm and 0.8cm]
        % Nodes
        \node (data) [box, align=center] {Carga de\\Datasets};
        \node (mask) [box, below of=data, align=center] {Máscara de Degradación\\(Random vs Nearby)};
        \node (opt) [box, right=of mask, align=center, fill=usmgold!20, draw=usmgold!80!black] {Optimización Optuna\\(18 Algoritmos)};
        \node (buffer) [box, right=of opt, align=center] {Buffer 2 seg.\\(Prevención Leakage)};
        \node (eval) [box, below=0.8cm of opt, align=center] {Evaluación Multi-Métrica\\(MAE, DTW, LSD)};
        
        % Arrows
        \draw [arrow] (data) -- (mask);
        \draw [arrow] (mask) -- (opt);
        \draw [arrow] (opt) -- (buffer);
        \draw [arrow] (buffer.south) |- (eval.east);
        \draw [arrow] (opt) -- (eval);
    \end{tikzpicture}
    
    \vspace{0.2cm}
    \textbf{Datasets Utilizados (Evaluación Externa):}
    \begin{table}[]
        \centering
        \small
        \begin{tabular}{lccc}
            \toprule
            \textbf{Dataset} & \textbf{Canales} & \textbf{Densidad} & \textbf{Desafío Principal} \\
            \midrule
            PhysioNet EEG & 64 & Alta & Referencia / Ideal \\
            MNE Sample & 60 & Alta & ERPs clínicos (N100, P300) \\
            BCI IV 2a & 22 & Baja & \textbf{Degradación Severa} \\
            \bottomrule
        \end{tabular}
    \end{table}
    
    \begin{itemize}
        \item \textbf{Fair Benchmarking:} \textit{Todos} los métodos (incluidas las Splines) fueron optimizados en hiperparámetros.
    \end{itemize}
    \note{Aquí ilustramos el flujo. No evaluamos algoritmos con 'parámetros por defecto'. Optimizamos todo el abanico, desde Splines hasta TRSS, usando búsqueda Bayesiana multi-objetivo en la métrica MAE y LSD. Luego, un buffer estricto separa el set de entrenamiento del de prueba.}
\end{frame}

\begin{frame}{Validación Temporal vs. K-Folds Espaciales}
    \textbf{¿Por qué validación temporal con Buffer en lugar de K-Folds aleatorio?}
    \vspace{0.2cm}
    \begin{itemize}
        \item Si hubiéramos hecho \textit{K-Folds aleatorio} a nivel de muestras, la autocorrelación de banda pasante del EEG habría causado \textbf{Data Leakage} masivo.
        \item Si hubiéramos hecho \textit{K-Folds por sujeto}, cambiaría la topografía de la malla por la anatomía individual, rompiendo la optimización del Grafo.
        \item \textbf{Nuestra Solución (Hold-Out Cronológico):} 
        \begin{enumerate}
            \item Optimizamos hiperparámetros ($\alpha, \beta$) en el 70\% inicial del registro.
            \item Excluimos 2 segundos (\textit{buffer gap}) para romper la autocorrelación temporal.
            \item Evaluamos en el 30\% final.
        \end{enumerate}
    \end{itemize}
    \note{Este es un argumento crucial para cualquier evaluador de Machine Learning. El k-fold aleatorio a nivel de muestras infla los resultados. El buffer de 2 segundos garantiza que la evaluación sea rigurosamente honesta, simulando el uso clínico: calibro ahora, predigo minutos después.}
\end{frame}

% --- SECCIÓN 4: RESULTADOS ---
\section{Resultados}

\begin{frame}{Resultados Globales de Rendimiento (Optimizados)}
    Incluso tras dar a todos los métodos sus \textbf{mejores parámetros teóricos}:
    
    \vspace{0.5cm}
    \centering
    \begin{tabular}{lcccccc}
    \toprule
    \textbf{Familia} & \textbf{MAE $\downarrow$} & \textbf{RMSE $\downarrow$} & \textbf{DTW $\downarrow$} & \textbf{SNR $\uparrow$} & \textbf{LSD $\downarrow$} \\
    \midrule
    Baseline Geom. & 0.050 & 0.140 & 0.87 & 19.3 dB & 2.25 \\
    GSP Espacial   & 0.086 & 0.208 & 1.22 & 14.1 dB & 2.23 \\
    \textbf{Temporal (TRSS)} & \textbf{0.038} & \textbf{0.104} & \textbf{0.74} & \textbf{20.8 dB} & \textbf{2.22} \\
    \bottomrule
    \end{tabular}
    
    \vspace{0.3cm}
    \begin{itemize}
        \item La ventaja de TRSS es \textbf{estructural}. Las Spherical Splines no logran acortar la brecha en fidelidad de amplitud ni de forma (DTW).
    \end{itemize}
    \note{Estos números prueban que la falta de inercia es una barrera estructural. Las splines, aunque perfectamente afinadas por Optuna, siguen cometiendo más errores de forma (DTW) y de amplitud (MAE) que la regularización temporal.}
\end{frame}

\begin{frame}{Frontera de Pareto: Amplitud vs. Forma (BCI IV 2a)}
    \begin{columns}
        \begin{column}{0.4\textwidth}
            \textbf{Precisión (MAE) vs. Morfología (DTW)}
            \begin{itemize}
                \item Muestra el compromiso inherente de los métodos clásicos.
                \item \textbf{Cluster Geométrico (Morado/Azul):} Forma una barrera alejada del origen. Si ajustas para mejorar MAE, arruinas el DTW.
                \item \textbf{Dominancia TRSS:} Se ubica en la esquina inferior izquierda. Rompe la dicotomía logrando mínimo error de amplitud \textit{simultáneamente} con mínima distorsión temporal.
            \end{itemize}
        \end{column}
        \begin{column}{0.6\textwidth}
            \centering
            \includegraphics[width=\linewidth, height=0.75\textheight, keepaspectratio]{../../../results_optuna_final/tradeoff_scatter_bci_iv.png}
        \end{column}
    \end{columns}
    \note{Este gráfico es contundente. La frontera de pareto muestra que los métodos geométricos tienen un muro: no pueden mejorar el MAE sin arruinar el DTW. TRSS se salta ese muro por completo gracias a la dimensión temporal.}
\end{frame}

\begin{frame}{Panorámica de Robustez: Mapa de Calor (BCI IV 2a)}
    \begin{columns}
        \begin{column}{0.35\textwidth}
            \textbf{Estabilidad en todas las condiciones}
            \begin{itemize}
                \item Columnas: \% de pérdida y topología (random/nearby).
                \item El colapso geométrico es visible en las columnas \textit{nearby} de alta pérdida (tonos claros = alto error).
                \item \textbf{TRSS} se mantiene con la franja más oscura (menor error) de forma consistente, demostrando ser inmune a la densidad de la pérdida.
            \end{itemize}
        \end{column}
        \begin{column}{0.65\textwidth}
            \centering
            \includegraphics[width=\linewidth, height=0.75\textheight, keepaspectratio]{../../../results_optuna_final/heatmap_performance_bci_iv.png}
        \end{column}
    \end{columns}
    \note{Para resumir el rendimiento general, este mapa de calor muestra la matriz de métodos vs escenarios. En escenarios fáciles, todos son oscuros. Pero al movernos a la derecha (40 porciento nearby), el resto se aclara (falla), mientras que TRSS retiene la precisión.}
\end{frame}

\begin{frame}{Reconstrucción de 1 Canal: Baja vs Alta Degradación (Escala Fija)}
    \vspace{-0.2cm}
    \textbf{Reconstrucción de 1 Canal Faltante} (Ground Truth en negro punteado). Escalas Y fijadas para visualización justa.
    
    \vspace{0.1cm}
    \begin{columns}[T]
        \begin{column}{0.5\textwidth}
            \centering
            \textbf{Pérdida Baja/Media (10\% Nearby)}
            \includegraphics[height=0.7\textheight, width=\linewidth, keepaspectratio]{../../../results_optuna_final/erp_timeseries_mne_sample_nearby_0.1.png}
        \end{column}
        \begin{column}{0.5\textwidth}
            \centering
            \textbf{Pérdida Severa (40\% Nearby)}
            \includegraphics[height=0.7\textheight, width=\linewidth, keepaspectratio]{../../../results_optuna_final/erp_timeseries_mne_sample_nearby_0.4.png}
        \end{column}
    \end{columns}
    \note{Aquí vemos a la izquierda un caso clínico común (10 porciento). MNE y Spline lo hacen bien porque hay muchos canales vecinos sanos para compensar, aunque ICA y Spline a veces lanzan picos fuera de escala. Pero a la derecha, en 40 porciento (un colapso de parche masivo), las Splines INVIERTEN la fase destruyendo el potencial. TRSS (línea azul) se aferra a la inercia y sigue la trayectoria correcta. Nota: MNE interpolate_bads usa Splines por debajo con parámetros por defecto, por eso rinde similar a la Spline morada.}
\end{frame}

\begin{frame}{Degradación bajo Pérdida Extrema (Colapso Geométrico)}
    \begin{columns}
        \begin{column}{0.4\textwidth}
            \textbf{Dataset BCI IV 2a (Solo 22 canales)}
            \vspace{0.3cm}
            \begin{itemize}
                \item Eje Y: Decibeles (SNR).
                \item Al alcanzar un 20-30\% de pérdida contigua en un casco de baja densidad, los baselines geométricos sufren un colapso exponencial.
                \item \textbf{TRSS mantiene resiliencia lineal} al inyectar información del pasado y futuro.
            \end{itemize}
        \end{column}
        \begin{column}{0.6\textwidth}
            \includegraphics[width=\linewidth,height=0.75\textheight,keepaspectratio]{../../../results_optuna_final/degradation_bci_iv_nearby_ratio.png}
        \end{column}
    \end{columns}
    \note{Este gráfico refuerza lo anterior, pero estadísticamente sobre todo el dataset. En 22 canales, perder el 40 porciento es perder un hemisferio entero. La spline se cae al abismo negativo. TRSS aguanta sobre los 10 dB gracias a la inercia temporal.}
\end{frame}

\begin{frame}{Estudio de Ablación: ¿Qué aporta exactamente el tiempo ($\beta$)?}
    Para aislar la contribución de la inercia temporal, se contrastó \textbf{TRSS-Completo ($\alpha, \beta$)} contra \textbf{TRSS-Espacial ($\alpha, \beta=0$)}.
    
    \vspace{0.3cm}
    \vspace{0.2cm}
    \begin{columns}[T]
        \begin{column}{0.5\textwidth}
            \begin{block}{Dominio Temporal (MAE)}
                \textbf{Mejora Modesta ($\sim 0.6\%$)} \\
                La componente temporal no está diseñada para reducir el error absoluto medio; la amplitud es gobernada casi por completo por el Laplaciano espacial ($\alpha$).
            \end{block}
        \end{column}
        \begin{column}{0.5\textwidth}
            \begin{block}{Dominio Espectral (LSD)}
                \textbf{Mejora Pronunciada ($+8.5\%$)}\\
                Aislar el tiempo rescata la \textit{Log Spectral Distance}. Evita ruido de alta frecuencia y saltos artificiales. \textbf{Tamaño de efecto grande} ($\delta = -0.76$, $p<0.005$).
            \end{block}
        \end{column}
    \end{columns}
    
    \vspace{0.3cm}
    \centering
    \begin{tabular}{lcc}
        \toprule
        \textbf{Variante} & \textbf{LSD Promedio $\downarrow$} & \textbf{DTW Promedio $\downarrow$} \\
        \midrule
        TRSS-Espacial ($\beta=0$) & 2.45 & 0.81 \\
        \textbf{TRSS-Completo ($\beta > 0$)} & \textbf{2.22} & \textbf{0.74} \\
        \bottomrule
    \end{tabular}
    \vspace{0.3cm}
    \textbf{Reflexión:} La regularización temporal asegura \textbf{comportamiento oscilatorio fisiológico} en altas frecuencias (bandas beta, gamma), evitando el \textit{jitter} artificial.
    \note{No exageramos nuestros resultados. Admitimos que el parámetro Beta mejora poco el error medio (MAE). Pero su función vital no es esa; su función es salvar la fase y las bandas de frecuencia altas (LSD). Esto es honestidad intelectual de altísimo nivel y le dará muchísima credibilidad a tu defensa.}
\end{frame}

% --- SECCIÓN 5: CONCLUSIONES ---
\section{Conclusiones y Trabajo Futuro}

\begin{frame}{Conclusiones}
    \begin{enumerate}
        \item \textbf{Validación de Hipótesis:} La regularización combinada de Grafo (topología) y Sobolev (inercia temporal) supera de forma robusta y estructural a la extrapolación euclidiana clásica.
        \item \textbf{Fair Benchmarking:} La optimización multi-objetivo cerró la brecha en escenarios benignos, pero evidenció el \textbf{colapso estructural} geométrico ante fallas contiguas severas.
        \item \textbf{Preservación Diagnóstica:} El beneficio real de la componente temporal reside en la preservación de las bandas espectrales (PSD) y la morfología transitoria (ERPs), cruciales para clínica y BCI.
    \end{enumerate}
    \note{Se confirman los objetivos: GSP y Tiempo superan a la esfera. Se hizo justamente. Y la relevancia clínica radica en la fase y frecuencia.}
\end{frame}

\begin{frame}{Trabajo Futuro}
    \begin{itemize}
        \item \textbf{Procesamiento BCI \textit{Online}:} TRSS procesa ventanas matriciales completas iterativamente. A futuro se deben implementar \textit{ventanas deslizantes causales} y portar el álgebra a \textbf{GPU (CuPy)}.
        \item \textbf{Validación Descendente (\textit{Downstream Tasks}):} Cuantificar empíricamente cómo la mejora espectral (LSD) de TRSS incrementa la exactitud (\%) en clasificadores profundos (ej. EEGNet).
        \item \textbf{Grafos Dinámicos y Deep Learning:} Investigar cómo Grafos Neuronales (GNNs) podrían aprender topologías dependientes del tiempo, reemplazando las constantes estáticas $\alpha, \beta$.
    \end{itemize}
    \note{Queda mucho por hacer: acelerar el cálculo, usar Deep Learning (EEGNet) para probar que clasificamos mejor los trials recuperados, y automatizar el grafo con GNNs.}
\end{frame}

\begin{frame}[plain]
    \begin{center}
        \Huge \textbf{Gracias por su atención} \\[1em]
        \large Universidad Técnica Federico Santa María \\
        Departamento de Ingeniería Electrónica \\[2em]
        \Large ¿Preguntas?
    \end{center}
    \note{Muchas gracias por escuchar, quedo a la espera de las preguntas del comité.}
\end{frame}

% --- BACKUP SLIDES ---
\appendix
\begin{frame}[plain, noframenumbering]{Backup: Optimización de Hiperparámetros}
    \begin{center}
    \begin{tabular}{lll}
    \toprule
    \textbf{Método} & \textbf{Parámetro} & \textbf{Rango (Log / Lineal)} \\
    \midrule
    Spherical Spline & \texttt{eps} (Suavidad Legendre) & $[10^{-8}, 10^{-2}]$ \\
    \midrule
    GSP Tikhonov & $\alpha$ (Regularización Espacial) & $[10^{-4}, 10^{2}]$ \\
    \midrule
    \multirow{2}{*}{TRSS} & $\alpha$ (Espacio), $\beta$ (Tiempo) & $[10^{-3}, 10]$, $[10^{-4}, 1]$ \\
                          & \texttt{n\_iter}, \texttt{lr} (Descenso) & $[50, 500]$, $[10^{-4}, 0.1]$ \\
    \midrule
    Grafo (KNNG) & $k$, $\sigma$ (Ancho kernel gaussiano) & $[2, N/2]$, $[0.01, 5]$ \\
    \bottomrule
    \end{tabular}
    \end{center}
    
    \vspace{0.3cm}
    Se generaron 30 trials por configuración usando TPE para TRSS/Tikhonov, y 10 trials usando NSGA-II para los métodos baselines.
\end{frame}

\begin{frame}[plain, noframenumbering]{Backup: Derivación del Gradiente de TRSS}
    La función objetivo de TRSS:
    \[ J(\mathbf{Z}) = \|\mathbf{S}\mathbf{X} - \mathbf{S}\mathbf{Z}\|_F^2 + \alpha \text{tr}(\mathbf{Z}^\top \mathbf{L} \mathbf{Z}) + \beta \|\mathbf{Z}\mathbf{D}_t\|_F^2 \]
    
    El gradiente respecto a $\mathbf{Z}$ es:
    \[ \nabla_{\mathbf{Z}} J = 2 \mathbf{S}^\top \mathbf{S} (\mathbf{Z} - \mathbf{X}) + 2 \alpha \mathbf{L} \mathbf{Z} + 2 \beta \mathbf{Z} \mathbf{D}_t \mathbf{D}_t^\top \]
    
    \begin{itemize}
        \item El operador $\mathbf{D}_t \mathbf{D}_t^\top$ corresponde al \textbf{Laplaciano de la cadena temporal} (grafos de ruta unidimensional).
        \item La matriz \textbf{S} fija las entradas conocidas, obligando a los canales perdidos a absorber el gradiente iterativamente.
    \end{itemize}
\end{frame}

\end{document}
